% META: {"id":"TCOMPL-ECH-ME-002","chapitre":"TCOMPL-ECHANTILLONNAGE","type_objet":"methode","capacites_codes":["C2"],"methodes":["M2"],"mode_creation":"ex_nihilo","sources_inspiration":[],"status":"approved","origin":"needs_review","status_history":["needs_review","audited","accepted","approved"],"release_acceptance":"RELEASE_OWNER_BATCH_ACCEPTANCE","acceptance_closure_digest":"sha256:9b3ccf9a81c5520fbb7e03b7b2d3e7bf2057a02834b6d908bf3be5f49f3b553f","acceptance_receipt_digest":"sha256:4fad86f0282e0fa4e0419cca1ad6379ae7af5a280c04a4a90880f90a1c044242"}
% BEGIN-VERIFY
% from sympy import binomial as C
% assert C(4,1) + C(4,2) == C(5,2) == 10
% assert C(5,2) + C(5,3) == C(6,3) == 20
% END-VERIFY
\begin{fichemethode}{M2}{Déterminer des coefficients binomiaux avec le triangle de Pascal}

\quandUtiliser{%
  Utiliser cette méthode lorsque l'énoncé demande un coefficient binomial
  $\binom{n}{k}$ pour de petites valeurs de $n$, ou de construire le triangle de
  Pascal ligne par ligne.
}

\pasApas{%
  \begin{enumerate}
    \item Écrire la ligne $n=0$ : elle contient le seul nombre 1.
    \item Chaque ligne commence et finit par 1 : $\binom{n}{0} = \binom{n}{n} = 1$.
    \item Relation de Pascal : chaque coefficient interieur est la somme des DEUX
      nombres situes au-dessus :
      $\binom{n}{k} = \binom{n-1}{k-1} + \binom{n-1}{k}$.
    \item Construire les lignes une a une jusqu'a la ligne $n$ voulue, puis lire
      $\binom{n}{k}$ en position $k$ (en comptant depuis $k=0$).
    \item Utiliser la symetrie $\binom{n}{k} = \binom{n}{n-k}$ pour limiter les calculs.
  \end{enumerate}
}

\exempleRedige{%
  Déterminer $\binom{5}{2}$ avec le triangle de Pascal.

  \medskip
  \commentaireMarge{Chaque nombre = somme des deux au-dessus.}%
  Lignes successives :
  $n=3$ : $1\ 3\ 3\ 1$ ;\quad
  $n=4$ : $1\ 4\ 6\ 4\ 1$ ;\quad
  $n=5$ : $1\ 5\ 10\ 10\ 5\ 1$.

  \medskip
  En position $k=2$ de la ligne $n=5$ : $\binom{5}{2} = 10$. On retrouve bien
  $\binom{4}{1} + \binom{4}{2} = 4 + 6 = 10$.
}

\pieges{%
  \begin{itemize}
    \item \textbf{Piège 1.} Decaler l'indice $k$ : la ligne commence a $k = 0$,
      donc $\binom{5}{2}$ est le TROISIEME nombre de la ligne 5.
    \item \textbf{Piège 2.} Écrire la relation de Pascal avec de mauvais indices :
      les deux termes ont pour indice superieur $n-1$, jamais $n$.
    \item \textbf{Piège 3.} Oublier les 1 des extrémités en construisant une ligne.
  \end{itemize}
}

\verifier{%
  Contrôler que la somme de la ligne $n$ vaut $2^n$ (ligne 5 : $1+5+10+10+5+1 = 32$)
  et que la ligne est symetrique.
}

\sEntrainer{\refExos{M2}}

\end{fichemethode}
